\section{Discussion}
\label{sec:discussion}

% TODO: Add findings discussion after analysis is complete.

\subsection{PT Looping: A Related Strategy Variant}

During whale identification, several addresses were found executing a structurally related but analytically distinct strategy using Pendle Principal Tokens (PT-sUSDe) as collateral rather than sUSDe directly. These are documented here for completeness but excluded from the core quantitative analysis.

\textbf{How Pendle works.} Pendle Finance splits a yield-bearing token into two components: a Principal Token (PT) redeemable for the underlying at a fixed maturity date, and a Yield Token (YT) that captures floating yield until maturity. Because the PT has been stripped of its yield, it trades at a discount implying a fixed annualised return for holders who hold to maturity.

\textbf{PT looping mechanism.} PT loopers use PT-sUSDe as Morpho collateral: obtain sUSDe, split via Pendle Router to receive PT-sUSDe and YT-sUSDe, supply PT-sUSDe as collateral on Morpho Blue, borrow stablecoin against it, convert borrowed stablecoin back to sUSDe, and repeat. The carry spread is the PT implied fixed rate minus the borrow rate, locked in at entry rather than floating daily.

\textbf{Why PT looping is excluded.}

\begin{table}[h!]
\centering
\caption{Comparison of YBT looping versus PT looping.}
\label{tab:pt-comparison}
\begin{tabular}{lp{4.5cm}p{4.5cm}}
\hline
\textbf{Dimension} & \textbf{YBT Looping (this study)} & \textbf{PT Looping (excluded)} \\
\hline
Collateral & Floating-rate YBT & Fixed-rate PT \\
Carry spread & Changes daily & Locked at entry \\
Price trajectory & Varies with yield & Predetermined accretion to par \\
Liquidation driver & YBT depeg or rate spike & PT price deviation from fair value \\
Panel structure & Continuous & Breaks at maturity \\
\hline
\end{tabular}
\end{table}

The maturity breaks are the primary engineering challenge. Each PT expiry requires special panel handling that does not apply to floating-rate positions. Combined with the need for Pendle AMM implied rate data (a separate API), this represents meaningful scope expansion.

Two PT loopers were identified during screening: an institutional-scale operator (\texttt{0x1b82b4f4...}) moving \$28--34M per atomic transaction with only 6 total transactions, and an iterative looper (\texttt{0xbfb77183...}) with 34 transactions concentrated in a two-day burst. Both are excluded from the core sample but demonstrate that recursive leverage extends beyond floating-rate YBTs into fixed-rate Pendle derivatives.

PT looping occupies a different position in the risk-return space: operators trade floating carry exposure for a predetermined fixed return. This places PT looping closer to fixed-income arbitrage than to the yield farming strategies analysed here, and represents a natural direction for future work.

\subsection{Limitations}

The most important limitation of the present analysis is the small number
of whales. With only seven wallets, several archetypes are dominated by a
single whale, most notably the extreme-leverage archetype, which is almost
entirely one wallet (W-2). Archetype and wallet identity are therefore
partially confounded, and the precise risk ordering between archetypes
should be regarded as provisional rather than settled. A larger and more
diverse set of wallets would be needed to establish whether the archetypes
generalise beyond individual operators.

This sample limitation also affects the stress comparison of
Section~\ref{sec:stress-by-archetype}. Because the supervised
\texttt{stressed\_state} label is defined relative to each whale's own
health-factor distribution, the stressed share of an archetype depends on
its wallet composition; this is why the per-whale label understates the
extreme-leverage archetype, which has no calm baseline of its own to stand
out against. We mitigate this by reporting an absolute health-factor
cut-off applied identically to every whale, but the cut-off itself is a
hand-picked buffer rather than one calibrated to observed liquidation
events. With more whales we would expect the two measures to converge, and
calibrating the absolute threshold to realised liquidations is left for
future work.

Further limitations shared with the supervised analysis include the
approximation of loop depth from observable position data and the
restriction to two lending protocols over a roughly seventeen-month window,
which together bound how far the archetype taxonomy can be generalised.
